\section[范畴化]{范畴化（5 月 27 日）}

\srcpage{81}

今天我们要讲范畴化（categorification）。它其实不是一个精确、良定义的概念：它是一族想法，与去范畴化（decategorification）互为反向，而后者是良定义的。

给你一个数，你可能会问：它是否来自某个带有更多结构的对象？例如，你算出的那个量，是否其实是某个向量空间的维数？当然，我们知道哪些数是维数；但我们想要的是一个在此语境中有趣的向量空间。可能的选择只有正整数——或者取 $\Z$-分次向量空间的次数维数或欧拉示性数，从而允许负整数。用这种方式从结构里取出数，就是去范畴化；反过来走，就是范畴化的一个例子。

范畴化的目的是添加有用的结构。例如，若你的向量空间上有算子 $U$ 作用，我们可以去范畴化并取迹，得到一个数。先验地，这个数与空间的维数毫无关系；但一旦范畴化，两者就有关系。同样，把一个数写成其他数之和并不特别；而把一个向量空间写成向量空间的直和却相当有用。你可能感兴趣的类似事情，是把群环的特征标范畴化为真正的表示。

范畴化有时能揭示更深层结构的影子：如果一个数意外地是正整数，也许它是某个尚未现身的向量空间的维数。

\srcpage{82}

再做一次这件事是有益的：把向量空间范畴化为范畴（通常带额外结构）。为理解这一点，我们需要去范畴化的概念，而 Grothendieck 群 $K_0$ 可以胜任。例如，若有类函数代数 $\C[G/G]$，你可能想把它范畴化为其表示范畴，并且 $\C[G/G]\cong K_0(\Rep_G)$。如果你能把一个向量空间实现为某个范畴的 $K_0$，就能得到额外结构，或许包括由范畴中单对象或投射对象给出的典范基。一个常见的例子，是把函数构成的向量空间范畴化为层范畴。

拓扑场论对范畴化有很多话要说，反之亦然。给定 TFT $Z$，在 $S^1$ 上作维数约化会得到另一个 TFT $Z'$，定义为 $Z'(X):=Z(S^1\times X)$。想法是：这个过程是去范畴化的一般形式：$Z'(X)$ 是 $Z(X)$ 的去范畴化。例如，若 $X$ 是闭 $(n-1)$-维流形，则 $Z'(X)$ 是复数，是 $Z(X)$ 的维数；特别地，它是正整数。

更一般地，若 $\cC$ 是对称幺半范畴，$Z\colon\mathrm{Bord}_n\to\cC$ 是 TFT，则可以问 $Z(M)$ 与 $Z(S^1\times M)$ 的关系，后者是 $\End_\cC(\mathbf{1})$ 的元素。公式相同：把 $S^1$ 分解为两段区间，先余赋值再赋值。

\begin{definition}
设 $\cC$ 是对称幺半范畴，$V\in\cC$。若存在对象 $V^\vee\in\cC$ 与映射 $c\colon\mathbf{1}\to V\otimes V^\vee$、$e\colon V^\vee\otimes V\to\mathbf{1}$（分别称为余赋值与赋值），使得 Zorro 之标记（mark of Zorro）公理被满足，则称 $V$ 可对偶化（dualizable）：（TODO：图）\yizhu{Zorro 之标记指对偶对象的两条 zigzag 恒等式；名字来自 Zorro 拔剑划出的 Z 字形。}{B（nLab “dualizable object”）}
\end{definition}

若 $V$ 可对偶化，则它有维数，定义为 $e\circ c\in\End_\cC(\mathbf{1})$（不一定是数！）。这在任何对称幺半范畴中都有意义。

\begin{example}
对 $\cC=\Vect_k$，$V^\vee$ 是通常的对偶：映入 $k$ 的映射。赋值是自然配对 $V^\vee\otimes V\to k$，由 $(\varphi,v)\mapsto\varphi(v)$ 给出。余赋值使用典范同构 $V\otimes V^\vee\cong\End(V)$，并把 $1$ 送到 $\id_V$。这关键地用到有限维性：你可以用一组基与对偶基写出单位矩阵。特别地，$\dim V=\Tr(\id_V)$：在同构 $\End(V)\cong V\otimes V^\vee$ 下，赋值 $\End(V)\to k$ 就是迹。
\end{example}

受此启发，我们把一般的迹映射 $V\otimes V^\vee=\End(V)\to\mathbf{1}$ 定义为赋值。

为了思考向量空间的维数，这用了不少范畴论；但要点是：我们对 TFT 中可能的取值了解了不少，也了解了去范畴化如何产生。

\begin{example}
设 $\cC$ 是 $k$ 上链复形的对称幺半范畴。则 $\cC$ 中维数的概念与欧拉示性数一致。
\end{example}

\begin{example}
让我们把事情升级一点：$k$-代数属于对称幺半 $2$-范畴 $\mathrm{Alg}_k$，其对象是 $k$-代数，态射是双模，$2$-态射是双模同态。双模的复合是张量积。在这个设定中，维数是什么？

首先，每个 $k$-代数 $A$ 都可对偶化：对偶是 $A^{\mathrm{op}}$。我们需要赋值与余赋值映射，它们分别是 $(k,A\otimes A^{\mathrm{op}})$-双模与 $(A\otimes A^{\mathrm{op}},k)$-双模。这两者都等价于一个左 $A$-作用与一个右 $A$-作用的数据，而这样的双模有自然选择，即 $A$ 本身。你可以验证它满足 Zorro 条件，所以在此意义下问 $A$ 的维数是有意义的。这将是 $\End_{\mathrm{Alg}_k}(k)$ 的元素，即 $(k,k)$-双模范畴（也就是向量空间）。具体地，先余赋值再赋值，得到
\[
A\otimes_{A\otimes_k A^{\mathrm{op}}}A,
\]
这就是所谓的 Hochschild 同调 $HH_*(A)$。有自然映射 $\tau\colon A\to HH_*(A)$，由 $a\mapsto a\otimes 1$ 给出，且 $\tau(ab)=\tau(ba)$，所以 $\tau$ 是迹。事实上，Hochschild 同调有泛性质：它是泛迹映射：任何迹映射 $A\to B$（即通过 $a,b\mapsto ab$ 与 $a,b\mapsto ba$ 的余等化子分解的映射）都通过 $HH_*(A)$ 与 $\tau$ 分解。

如果你在非导出设定中思考这件事，你只看到 $HH_0$，它就是 $A$ 的余中心（cocenter）：中心是满足 $ab=ba$ 的泛「映入」对象，而余中心是自然的「射出」对象。
\end{example}

所以 $HH_*$ 是去范畴化的一种形式。确切地说，我们在去范畴化 $\mathrm{Mod}_A$：存在 $2$-范畴的嵌入 $\mathrm{Alg}_k\to\mathrm{Cat}_k$，把 $A$ 送到 $\mathrm{Mod}_A$。于是对模范畴我们有了维数的概念（事实上，可以直接对 $\mathrm{Cat}_k$ 中可对偶化的范畴定义 Hochschild 同调）。但模范畴很好，所以有这个视角是好事。

\srcpage{83}

现在我们已经说出两种去范畴化范畴的方式：$K$-理论 $K_*$ 与 Hochschild 同调 $HH_*$。后者更弱，却容易计算得多。$K$-理论，尤其是高阶 $K$-理论，不是线性对象，可能极难计算。在表示论中遇到的各类范畴里，$HH_*(\cC)\cong k\otimes_\Z K_*(\cC)$；例如，对有限群的表示范畴，这是对的（TODO：至少在 $K_0$ 上）。

喜欢 $HH_*$ 的另一个重要理由（尽管它比 $K$-理论弱）是：Hochschild 同调在场论中有归宿，而 $K$-理论没有。

使用 $\infty$-范畴的一个好处是：我们总可以忘掉不可逆的态射，所以记录范畴的阶数没那么重要：在任意范畴层级，都可以用一些相同的手段定义对象的维数。这在非导出设定中完全不成立，那里复合的非结合性可能是大得多的问题。这对推广向量空间的部分故事也有用：如果理论建立得当，$A\otimes_k A^{\mathrm{op}}\cong\End(\mathrm{Mod}_A)$。因此从自函子到 $k$ 的赋值映射配得上「迹」之名；它吐出一个向量空间，而 $\Tr(\id_{\mathrm{Mod}_A})=HH_*(A)$。经典地、更具体地说，这意味着对任意 $(A,A)$-双模 $M$，我们附上迹
\[
\Tr(M):=M\otimes_{A\otimes A^{\mathrm{op}}}A,
\]
通常记作 $HH_*(A,M)$；于是 $\dim A=HH_*(A,A)=\Tr(\id_{\mathrm{Mod}_A})$。

回到场论：我们纯粹抽象地、只用对偶数据定义了维数的概念。顾名思义，这一数据被对称幺半函子保持，因此 $Z(\dim(X))=\dim(Z(X))$。它保证存在，因为 $\mathrm{Bord}_n$ 中所有对象都可对偶化！对偶是取相反定向的流形，赋值与余赋值映射是柱体，边界被看作纯入射或纯出射，从而得到 $\varnothing$ 与 $M\sqcup(-M)$ 之间的映射。$M$ 在 $\mathrm{Bord}_n$ 中的维数是 $S^1\times M$。只要 $M$ 不是最高维，你可以在任意范畴层级做这件事（这样我们才能用配边构造出 $M\times S^1$）。

许多情形下，我们的场论 $Z$ 通过一个函子分解：它取场的空间，取值于空间的对应范畴。再通过取函数、层等线性化到目标 $\cC$，例如把 $M$ 送到 $\LocSys_{G^\vee}(M)$，再送到它上面的拟凝聚层 $\QCoh(\LocSys_{G^\vee}(M))$。这两个函子都是对称幺半的，于是可以问维数在它们之下是什么样子。

但在此之前，先谈 $\mathrm{Bord}_n$ 中的迹。你取自映射的迹，即配边 $M\to M$。例如，可以通过取微分同胚的映射柱得到。无论如何，给定配边 $X\colon M\to M$，我们可以把两份边界 $M$ 等同起来，得到闭 $n$-维流形；这就是 $X$ 的迹。若 $X$ 是微分同胚的映射柱，其迹就是映射环面（mapping torus）。

在对应的世界里也有迹，有时叫做非线性空间。设 $X\leftarrow Z\to Y$ 是几何对象（空间？簇？叠？意象的对象？此刻确切的细节并不重要）的某个图表。这个对应范畴中任何 $X$ 不仅可对偶化，而且自对偶：即指定了同构 $X\cong X^\vee$。这是因为对角给出对应 $\mathrm{pt}\leftarrow X\xrightarrow{\Delta}X\times X$，它就是余赋值，而赋值类似。

于是 $X$ 可对偶化；它的维数是什么？这是一个点的自对应。对应的复合是拉回，所以我们最终得到拉回方块
\[
\begin{tikzcd}
X\times_{X\times X}X \arrow[r] \arrow[d] & X \arrow[d, "\Delta"] \\
X \arrow[r, "\Delta" ] & X\times X
\end{tikzcd}
\]
再把两条投影经 $X$ 与 $X\times X$ 送到 $\mathrm{pt}$，就得到点的自对应。这个拉回公式看起来很像 Hochschild 同调的公式，这并非巧合；我们不必操心 $(-)^{\mathrm{op}}$，因为 $X$ 自对偶。我们这里得到的，在导出几何设定中解释，是对角线与自身在 $X\times X$ 中的相交。这叫做环路对象（loop object）或惯性对象（inertia object）$LX$。

为什么这应当是环路空间？若 $X$ 是 $M$ 上的场空间，通常形如 $\mathrm{Map}(M,??)$，则 $\dim X$ 是 $S^1\times M$ 上的场空间，即 $\mathrm{Map}(S^1\times M,??)$，也就是 $\mathrm{Map}(M,??)$ 上的环路。

\srcpage{84}

现在推广到对应中映射 $f\colon X\to X$ 的迹。这对应于一个对应 $X\leftarrow X\to X$。则迹是 $X|_{\Delta\hookrightarrow X\times X}$，即 $f$ 的图像与对角线的相交。这与迹的通常概念相合：给定有限集 $X$ 与映射 $f\colon X\to X$，$\Tr(f_*\colon k[X]\to k[X])$ 是对角元素之和，即不动点的个数。这也是 Lefschetz 不动点公式中出现迹的同一理由。于是在我们这三个世界里：$\mathrm{Bord}_n$ 中的迹是映射环面；$\mathrm{Corr}$ 中的迹是不动点；而 $\mathrm{Cat}$ 或 $\mathrm{Sh}$ 中的迹是迹或 Hochschild 同调。

线性化的一种办法是取层，定义对称幺半函子 $\mathrm{Corr}\to\mathrm{Cat}$：给定对应 $X\xleftarrow{p}Z\xrightarrow{q}Y$，映射 $\mathrm{Shv}(X)\to\mathrm{Shv}(Y)$ 是推拉 $q_*p^!$。我们想说这与取维数或取迹交换，但这需要一个高度非平凡的事实：
\[
\mathrm{Shv}(X\times X)\xrightarrow{\ ?\ }\mathrm{Shv}(X)\otimes\mathrm{Shv}(X).
\]
这如此非平凡，以至于它在若干重要设定中是错的，必须绕着走。我们这学期在 BZR 讨论班讨论的内容与此相关：人们花了很多功夫试图证明迹的猜想，而不是观察到迹的同义反复，很大程度就是因为这个。

几何 Satake 定理是原始 Satake 定理的范畴化。范畴化大致在 1990 年代被首次说出，有 Khovanov、Baez 以及许多其他人的工作；但它也有可追溯到 Grothendieck 的先声，以及诸如 Lefschetz 不动点公式这样的经典想法。这个公式说：你可以通过计算诱导映射在上同调上的迹，来数映射 $f\colon X\to X$ 的不动点个数。（如所述这并不完全正确，但要点是两者密切相关，并且可以通过添加一些形容词使之精确。）公式是 $H^i$ 上的迹乘以 $(-1)^i$——但若你在链复形的对称幺半范畴中理解迹，这个带符号的分次迹恰好就是你得到的，所以很好：这真的是上同调上的迹。

不动点的个数是计数 $\#(\Delta\cap\Gamma(f))$（或 $H^*(\Delta\cap\Gamma(f))$ 的维数），即 $f$ 在对应范畴中的迹。所以在某种意义上，Lefschetz 定理是取上同调的对称幺半性。

这是经典 Lefschetz 不动点公式。接下来是 Grothendieck–Lefschetz 不动点公式：设 $X$ 是 $\F_q$ 上的簇；则可以通过上同调上的迹确定 $X$ 的点数。具体地，
\[
\#X(\F_{q^r})=\Tr\bigl(\varphi^r\mid H^*_c(X;\Ql)\bigr).
\]
确切地说，我们取 Frobenius 的 $r$ 次幂的紧支 $\ell$-进 étale 上同调；但对我们要说的是：点数就是迹。某种意义上，$H^*_c(X)$ 连同 Frobenius 作用把点数范畴化为一个带自同态的向量空间，而我们通过取迹去范畴化。而紧支上同调是前推 $\pi_!\Ql$ 到一个点；这就是所谓函子—层对应（function-sheaf correspondence）的起源，它是范畴化较深刻的显现之一，来自 Weil 猜想。

要更深入函子—层对应，我们考虑 $\F_q$ 上概形（叠）上的函数：有了这个，就可以取任意 $r$ 的 $\F_{q^r}$-点。（紧支）函数有积分映射，产生一个数；而这些概形上的层有紧支上同调，给出带 Frobenius 自同态的向量空间，去范畴化为一个数。

函子—层对应要求填上方图的最后一条边：在层上取迹，产生函数，使得函数的积分等于我们所开始的层的上同调上 Frobenius 的迹：
\[
\Tr(\varphi,\cF)(x)=\Tr(\varphi,i_x^*\cF).
\]
为使这有意义，$\cF$ 必须是 Weil 层：它在 Frobenius $\varphi$ 下弱等变；且 $x\in X(\F_q)$。这类似于构造层（constructible sheaf）的世界：在每个茎上取欧拉示性数，就能在每一点得到一个数。于是单个层给你无穷多个函数：对所有 $r$ 在 $\F_{q^r}$-点上做这件事。而范畴化暗示：有趣的函数应当来自层，而且我们应当能得到一些仅在函数层面看不到的额外有趣结构。这个想法在表示论中真正起飞，是随着 Lusztig、Deligne–Lusztig 与 Kazhdan–Lusztig 的工作：一句口号式的要点是，若 $G$ 是约化群（例如 $\mathrm{GL}_n$ 或 $\mathrm{PSL}_n$），有限群 $G(\F_q)$（这是获得大量有趣有限群的一条途径）的表示论是「几何的」，意思是它有范畴化：它来自把 Frobenius 的迹作用于层。

\srcpage{85}

这一想法有种种显现，但一般原则是：这些群的表示论在陪集空间上的层中有几何起源。特别地，几乎所有有限单群都是某个 $q$ 与 $G$ 的 $G(\F_q)$——除了交错群（其表示论我们已经理解）与 sporadic 群。Lusztig 用这一点给出了这些群的特征标表的描述——这是相当重要而令人印象深刻的成就。

例如，Lusztig 的特征标层（character sheaf）表明：$\C$ 值 $G(\F_q)$-表示的特征标，是 $G/G$ 上关于共轭等变的层的迹；共轭等价意味着迹是类函数。这个对应并不精确：单对象未必相同；但它相当有力。

另一个例子是 Kazhdan–Lusztig 理论：有限 Hecke 代数是
\[
H_{G(\F_q),B(\F_q)}=\Ql[B(\F_q)\backslash G(\F_q)/B(\F_q)].
\]
他们把此实现为 $\Ql[G(\F_q),B(\F_q)]$ 的自同态，而 Hecke 代数的模出现为（TODO：？？）。要点是：我们得不到所有 $G(\F_q)$-表示，但能得到所有那些来自旗簇上函数的表示。Kazhdan–Lusztig 证明这个有限 Hecke 代数经由范畴化出现：它来自 Frobenius 在 $\mathrm{Shv}(B(\F_q)\backslash G(\F_q)/B(\F_q))$ 上的迹（这里 Frobenius 实际上是平凡的，所以这也是该层范畴的 $K_0$）。层范畴上的额外结构，例如 Verdier 理论，除其他东西外还给出 Hecke 代数的自然基：取范畴中的单对象。这个基性质相当好，并且知道 Hecke 代数表示论的不少信息。用 Hecke 范畴替换 Hecke 代数，是范畴化的一个关键例子，而且并非凭空而来；它是研究有限群表示论这一故事的一部分。

几何 Satake 定理是「Satake 对应是几何的」这一论断。展开它，我们曾有一个 Hecke 代数，叫做球面 Hecke 代数
\[
\mathcal{S}\mathrm{ph}:=\Ql[G(\cO)\backslash G(K)/G(\cO)];
\]
这个代数的结构基本不依赖于 $k$，只通过其剩余特征标而依赖；而 $k$ 是 $\F_q((t))$ 之类，或 $\Q_p$ 的有限扩域。要点是下面这个同构：
\[
\mathcal{S}\mathrm{ph}\cong\End_{G(K)}\mathrm{Fun}(G(K)/G(\cO)).
\]
因此球面 Hecke 代数作用于 $G(F)\backslash G(\A)/G(\cO_\A)$ 上的函数。

「几何 Satake 定理意味着 Satake 对应是几何的」这一想法是说：球面 Hecke 代数是球面 Hecke 范畴的迹，而后者是 $G(\cO)\backslash G(K)/G(\cO)$ 上层的范畴。这里 $K=\F_q((t))$，$\cO=\F_q[[t]]$——特征 $p$ 在这里要紧。混合特征情形也有一个版本，但困难得多，这归功于 Fargues–Scholze 的近期工作。

现在几何 Satake 把球面 Hecke 范畴等同于 $\Rep_{G^\vee}$。要去范畴化，我们取 Frobenius 的迹；Frobenius 是恒等，所以这就是 $K_0$。于是球面 Hecke 代数被等同于表示环 $K_0(\Rep_{G^\vee})$，这就是原来的 Satake 定理。你可以把它想成环面 $\cO(G^\vee/G^\vee)=\cO(T^\vee)^W$ 上的 Weyl 群不变函数。现在这告诉我们球面 Hecke 代数交换的另一个理由。

Langlands 对偶群的出现指示了某种深刻的东西：$\mathcal{S}\mathrm{ph}\cong K_0(\Rep_{G^\vee})$，其中 $G^\vee$ 由 $G$ 的对偶根数据组合地构造。这很有趣，因为 $K_0(\Rep_{G^\vee})$ 一般不能决定 $G^\vee$：群并不能由它的表示环决定。另一方面，几何 Satake 定理确实能恢复 $G^\vee$。这是范畴化的几何版本更强的几个理由之一。Lusztig 的名字又出现了：他是最早研究仿射 Grassmann 簇上反常层的人之一，而这是几何 Satake 对应故事的重要部分。
